%=======================================================================
% cap2_trabajo_relacionado.tex
%=======================================================================

Este capítulo revisa el estado del arte en las áreas que sustentan el
proyecto, avanzando desde lo general a lo específico, e identifica la brecha
de la literatura que ofrece la oportunidad de investigación. La revisión se
organiza en tres subáreas: identificación de aves por audio (bioacústica
computacional), identificación por imagen (incluida la fusión multimodal) y
modelos de distribución de especies (SDM). Cada sección describe brevemente la técnica,
revisa los trabajos más influyentes y destaca los aspectos que condicionan el
diseño de esta propuesta. El capítulo cierra con una tabla comparativa
(\tabref{tab:brecha}) y la formulación explícita de la brecha.

%-----------------------------------------------------------------------
\section{Identificación de aves por audio (bioacústica computacional)}
\label{sec:rel-audio}

La bioacústica computacional estudia la detección y clasificación automática
de sonidos animales a partir de grabaciones. Su desarrollo con aprendizaje
profundo ha seguido de cerca al de la clasificación general de audio: la
señal se transforma en una representación tiempo--frecuencia (habitualmente
un espectrograma en escala mel, \secref{sec:teoria-audio}), que se procesa como una imagen mediante redes
convolucionales (CNN, \secref{sec:teoria-cnn}) o, más recientemente, mediante
\textit{transformers} de audio (\secref{sec:teoria-transformers}).

\subsection*{Representaciones de entrada}

El espectrograma mel comprime el eje de frecuencia según una escala
perceptual y expresa la energía en escala logarítmica, lo que concentra la
información relevante para vocalizaciones animales y permite reutilizar
arquitecturas de visión por computador. En este proyecto se adopta el
espectrograma log-mel como representación principal por su compatibilidad
con los modelos preentrenados que se describen a continuación.

\subsection*{Preentrenamiento a gran escala y modelos de referencia}

La disponibilidad de AudioSet, un catálogo de más de dos millones de clips
de audio anotados con 527 clases de sonidos \citep{gemmeke2017audioset},
permitió demostrar que las arquitecturas CNN diseñadas para imágenes
se transfieren de forma efectiva a la clasificación de audio. Sobre esta base se construyeron familias de modelos
de audio preentrenados. Posteriormente, el \textit{Audio Spectrogram
Transformer} (AST) aplicó directamente un \textit{Vision Transformer} a
parches del espectrograma y alcanzó el estado del arte en AudioSet
\citep{gong2021ast}. Su arquitectura se describe en la
\secref{sec:teoria-transformers}. Estos modelos son candidatos naturales para
la identificación bioacústica de especies, dado que sus representaciones
internas codifican patrones espectrotemporales genéricos que pueden
especializarse con relativamente pocos datos.

En el dominio específico de las aves, BirdNET es el modelo más difundido:
una red convolucional profunda con conexiones residuales
(\secref{sec:teoria-cnn}) entrenada sobre espectrogramas de Xeno-canto y
otras colecciones, que en su versión original reconocía 984 especies de
Norteamérica y Europa en grabaciones de campo \citep{kahl2021birdnet}. Las
versiones posteriores superan las \num{6000} especies. Su adopción en
monitoreo acústico pasivo es masiva, pero su desempeño varía fuertemente
entre especies y contextos, por lo que debe validarse localmente antes de
usarse para inferencia ecológica. La
primera evaluación mundial de BirdNET, con \num{4224} grabaciones de 67
sitios en todos los continentes anotadas por expertos locales, confirmó esta
heterogeneidad (\tabref{tab:birdnet}): el desempeño es máximo en los
continentes de donde provienen la mayoría de sus datos de entrenamiento y
cae en Sudamérica \citep{funosas2026birdnet}. Estos resultados justifican
entrenar o adaptar modelos específicos para la avifauna chilena en lugar de
depender de un clasificador global.

\begin{table}[H]
  \centering\small
  \caption{Desempeño de BirdNET por región en la evaluación mundial de
  \citet{funosas2026birdnet}: área bajo la curva precisión--\textit{recall}
  (AUC-PR) por continente; el \textit{recall} varía entre 0,24 y 0,52 según
  el continente.}
  \label{tab:birdnet}
  \begin{tabular}{@{}l l l@{}}
    \toprule
    Región & AUC-PR & Cobertura en los datos de entrenamiento \\
    \midrule
    Norteamérica, Oceanía y Europa & 0,16--0,23 & alta \\
    Centro y Sudamérica & 0,13 & moderada \\
    África y Asia & inferior a 0,13 & baja \\
    \bottomrule
  \end{tabular}
\end{table}

\subsection*{\textit{Embeddings} bioacústicos transferibles}

Una línea reciente muestra que los clasificadores de aves a gran escala
producen \textit{embeddings} que se transfieren mejor a nuevas tareas
bioacústicas que los de modelos de audio general: entrenando clasificadores
lineales sobre las representaciones de BirdNET y de Perch (un modelo de
Google entrenado en Xeno-canto), Ghani \etal\ obtuvieron desempeño superior
a AudioMAE y otros modelos generativos en tareas de aves, anfibios,
mamíferos marinos e insectos, incluso en régimen de pocos ejemplos por clase
\citep{ghani2023embeddings}. El \textit{benchmark} BIRB formaliza esta
evaluación como un problema de recuperación de información con
desplazamiento de distribución entre grabaciones focales y paisajes sonoros
\citep{hamer2023birb}. Esta evidencia importa para las especies con pocas grabaciones, como las
amenazadas de Chile, y justifica considerar el \textit{fine-tuning} de
modelos preentrenados en aves como estrategia complementaria al AST.

\subsection*{Evaluación en condiciones de campo y \textit{benchmarks}}

Los desafíos anuales BirdCLEF son el \textit{benchmark} de referencia: sus ediciones recientes plantean la identificación de especies en
paisajes sonoros continuos de regiones con pocos datos (por ejemplo, África
oriental en 2023), donde los participantes deben lidiar con etiquetas
débiles, solapamiento de vocalizaciones y desplazamiento de dominio entre
las grabaciones focales de Xeno-canto usadas para entrenar y los paisajes
sonoros de prueba \citep{kahl2023birdclef}. Las soluciones ganadoras
combinan CNN preentrenadas, aumentación de datos y ensamblado de modelos.
Xeno-canto, la fuente principal de estos desafíos y de este proyecto, es una
colección de ciencia ciudadana con anotación débil (una especie focal por
grabación, con posibles especies de fondo), duración muy variable y fuerte
desbalance entre especies \citep{vellinga2015xenocanto}.

\subsection*{Desbalance y clases con pocos ejemplos}

El desbalance de clases es inherente a los datos bioacústicos: las especies
comunes acumulan miles de grabaciones y las raras apenas decenas. Entre las
estrategias establecidas está la aumentación en el dominio del
espectrograma, como SpecAugment, que enmascara bandas de tiempo y frecuencia
\citep{park2019specaugment}. Otra es la interpolación de ejemplos mediante
\textit{mixup} \citep{zhang2018mixup}. Una tercera son las funciones de
pérdida que reponderan las clases, como la pérdida focal
\citep{lin2017focal} o la pérdida balanceada por número efectivo de
muestras \citep{cui2019classbalanced}. Estas técnicas conectan directamente con el
hallazgo de escasez de audio en especies amenazadas que se reporta en el
\capref{cap:datos}.

%-----------------------------------------------------------------------
\section{Identificación de aves por imagen y fusión multimodal}
\label{sec:rel-imagen}

\subsection*{Clasificación \textit{fine-grained}}

La identificación de aves por fotografía es el caso paradigmático de la
clasificación visual \textit{fine-grained}: distinguir subcategorías de una
misma clase de objetos (especies dentro del grupo ``aves'') que comparten la
estructura global y difieren en detalles locales, como la forma del pico o
los patrones del plumaje, con alta variabilidad dentro de cada clase por
pose, iluminación y fondo. El conjunto
CUB-200-2011, con \num{11788} imágenes de 200 especies norteamericanas, ha
sido durante más de una década el \textit{benchmark} de referencia.
iNaturalist 2017, con más de \num{850000} imágenes de \num{5089} especies y
datos fuertemente desbalanceados entre clases, acercó el problema a las
condiciones reales de los datos de ciencia ciudadana
\citep{vanhorn2018inaturalist}. En ambos, el \textit{transfer learning}
desde modelos preentrenados en ImageNet es la práctica estándar. Las
arquitecturas residuales establecieron la base. EfficientNet introdujo el
escalado compuesto de profundidad, ancho y resolución, con mejor precisión y
menos parámetros \citep{tan2019efficientnet}. El \textit{Vision Transformer}
mostró que un \textit{transformer} aplicado a parches de imagen supera a las
CNN cuando se preentrena a gran escala \citep{dosovitskiy2021vit}. Las bases
de estas arquitecturas se presentan en el \capref{cap:teoria}.

\subsection*{Estrategias de fusión multimodal}

El aprendizaje multimodal busca representar y combinar información de
distintas modalidades. Según el momento en que ocurre esa combinación,
\citet{baltrusaitis2019multimodal} distinguen tres estrategias:
\textit{early}, \textit{intermediate} y \textit{late fusion}. La primera
concatena las características de entrada y entrena un único clasificador. La
segunda hace interactuar las representaciones latentes de cada rama antes de
decidir. La tercera entrena un clasificador por modalidad y combina sus
probabilidades. Esta última es robusta cuando falta una modalidad en
inferencia, porque basta usar la predicción del clasificador disponible,
permite entrenar cada rama con datos no emparejados y es simple de calibrar.
Su limitación es que no modela interacciones de bajo nivel entre las señales
de entrada. Las tres se formalizan en la \secref{sec:teoria-fusion} y se
comparan empíricamente en el \capref{cap:resultados}.

\subsection*{Fusión audio--imagen en fauna}

La literatura que combina audio e imagen para identificación de animales es
reciente y limitada. El \textit{benchmark} SSW60 (Sapsucker Woods 60), de
Cornell, es el primer conjunto que permite clasificar las mismas 60 especies
de aves en imagen, audio y video. Sus autores muestran que la fusión
audiovisual supera a cada modalidad por separado y que el conocimiento
se transfiere entre modalidades \citep{vanhorn2022ssw60}. Pero sus especies son
norteamericanas y sus datos audiovisuales provienen de video emparejado. No se identificaron trabajos que combinen audio e imagen para
avifauna sudamericana ni que lo hagan a partir de fuentes no emparejadas
(Xeno-canto e iNaturalist), como propone este proyecto.

%-----------------------------------------------------------------------
\section{Modelos de distribución de especies (SDM)}\label{sec:rel-sdm}

Los SDM relacionan registros de ocurrencia de una especie con variables
ambientales para estimar su distribución potencial en el espacio y,
eventualmente, en el tiempo. Constituyen una
herramienta central de la ecología y la planificación de la conservación.

\subsection*{MaxEnt y el problema de solo presencia}

La mayoría de los datos disponibles, incluidos los de ciencia ciudadana,
son de solo presencia: registran dónde se observó una especie, pero no dónde
se buscó sin éxito. MaxEnt aborda este problema con una idea simple: entre todas las formas de
repartir la probabilidad de presencia sobre el territorio, elige la más
``neutra'' posible (la de máxima entropía) que sea compatible con lo
observado, es decir, con que las condiciones ambientales promedio de los
sitios donde se vio la especie coincidan con las del modelo. El resultado es
una probabilidad de presencia que crece de forma exponencial con una
combinación de las covariables, formalmente equivalente a una regresión
logística o de Poisson regularizada \citep{phillips2006maxent}. La
\secref{sec:teoria-sdm} presenta sus ecuaciones. Un \textit{benchmark} reciente con
código reproducible confirma que MaxEnt bien configurado sigue estando entre
los métodos de mejor desempeño para datos de solo presencia, junto con
\textit{boosted regression trees} y ensambles \citep{valavi2022benchmark}, lo
que lo convierte en un \textit{baseline} exigente y no trivial.

\subsection*{Sesgo de muestreo y pseudoausencias}

Como los datos solo dicen dónde se vio a la especie, el modelo necesita
puntos de contraste, las \emph{pseudoausencias} o puntos de fondo, que
representan los lugares donde la especie podría no estar. Si esos puntos se
reparten al azar por el territorio aparece un problema: los registros de
ciencia ciudadana se concentran en zonas accesibles y pobladas, de modo que
el modelo aprende dónde hay observadores en lugar de dónde hay hábitat
adecuado. La solución más usada es elegir los puntos de fondo con el mismo
sesgo que las presencias, por ejemplo, a partir de registros de otras
especies del mismo grupo, lo que mejora sustancialmente las predicciones
\citep{phillips2009bias}. La \figref{fig:fondo-objetivo} del
\capref{cap:teoria} ilustra la diferencia. El número y la ubicación de las
pseudoausencias también influyen en el desempeño, y para eBird existen guías
que recomiendan filtrar por protocolo y esfuerzo y submuestrear espacialmente
\citep{johnston2021ebird}.

\subsection*{SDM con aprendizaje profundo}

Los SDM neuronales reemplazan las funciones de respuesta paramétricas de la
regresión (una curva fija por covariable) por redes que aprenden
interacciones no lineales entre covariables y que pueden compartir
representaciones entre muchas especies, como se formaliza en la
\secref{sec:teoria-sdm}. Botella \etal\ propusieron
el primer enfoque de este tipo a escala continental \citep{botella2018deepsdm}.
Deneu \etal\ demostraron que las CNN que reciben parches de imágenes
ambientales alrededor de cada ocurrencia superan a los modelos puntuales al
capturar la estructura espacial del paisaje \citep{deneu2021cnn}.
Estopinan \etal\ mostraron que incorporar series temporales de Sentinel-2
aporta información estacional que beneficia especialmente a las especies con
pocas ocurrencias \citep{estopinan2022sentinel}. En paralelo, las
representaciones neuronales implícitas espaciales (SINR) codifican
directamente coordenadas y fecha en una red multiespecie entrenada con
registros de iNaturalist, produciendo mapas globales de miles de especies
\citep{cole2023sinr}. Específicamente para aves, SatBird combina imágenes
satelitales y datos de eBird en un conjunto de referencia para predecir tasas
de encuentro de cientos de especies \citep{teng2024satbird}.

\subsection*{Validación espacial}

Los datos de ocurrencia presentan autocorrelación espacial: registros
cercanos comparten ambiente y, si se reparten al azar entre entrenamiento y
prueba, el modelo obtiene estimaciones de desempeño optimistas. Roberts
\etal\ argumentan que la validación cruzada debe respetar la estructura de
los datos mediante bloques espaciales, temporales o jerárquicos
\citep{roberts2017crossvalidation}. El paquete \texttt{blockCV} implementa
la generación de \textit{folds} espacialmente separados y la estimación del rango de
autocorrelación para dimensionar los bloques \citep{valavi2019blockcv}. Este
protocolo se adopta para la comparación entre el SDM neuronal y MaxEnt.

%-----------------------------------------------------------------------
\section{Brecha de la literatura}\label{sec:brecha}

La \tabref{tab:brecha} resume los trabajos más representativos revisados,
según modalidad, región geográfica, número de especies y tarea abordada.

\begin{table}[!htb]
  \centering\footnotesize
  \caption{Comparación del trabajo relacionado y posicionamiento de esta
  propuesta.}
  \label{tab:brecha}
  \begin{tabularx}{\textwidth}{@{}>{\raggedright\arraybackslash\hsize=1.25\hsize}X
    >{\raggedright\arraybackslash\hsize=0.9\hsize}X
    >{\raggedright\arraybackslash\hsize=0.85\hsize}X
    r
    >{\raggedright\arraybackslash}X@{}}
    \toprule
    Trabajo & Modalidad & Región & N.º esp. & Tarea \\
    \midrule
    \citet{kahl2021birdnet} (BirdNET) & Audio & Norteamérica, Europa & 984 & Clasificación \\
    \citet{ghani2023embeddings} & Audio & Global (Xeno-canto) & $>$\num{10000} & Transferencia de \textit{embeddings} \\
    \citet{kahl2023birdclef} (BirdCLEF) & Audio & África oriental & 264 & Detección en paisajes sonoros \\
    \citet{wah2011cub} (CUB-200) & Imagen & Norteamérica & 200 & Clasificación \textit{fine-grained} \\
    \citet{vanhorn2018inaturalist} & Imagen & Global & \num{5089} & Clasificación con datos desbalanceados \\
    \citet{vanhorn2022ssw60} (SSW60) & Audio, imagen, video & Norteamérica & 60 & Clasificación audiovisual \\
    \citet{phillips2006maxent} (MaxEnt) & Ambiental & General & n.a. & Distribución (solo presencia) \\
    \citet{deneu2021cnn} & Ambiental (parches) & Francia & miles (flora) & Distribución (SDM profundo) \\
    \citet{cole2023sinr} (SINR) & Coordenadas + fecha & Global & \num{47375} & Distribución multiespecie \\
    \citet{teng2024satbird} (SatBird) & Satelital + eBird & EE.\,UU./Kenia & 670 & Tasas de encuentro \\
    \midrule
    \textbf{Este trabajo} & \textbf{Audio, imagen, ambiental} & \textbf{Chile continental} & \textbf{55} & \textbf{Clasificación + distribución} \\
    \bottomrule
  \end{tabularx}
\end{table}

A partir de esta revisión se identifican tres brechas concretas. La primera
es geográfica: los clasificadores bioacústicos y visuales de referencia se
entrenan y evalúan sobre avifaunas de Norteamérica y Europa, y la evidencia
global muestra que su desempeño se degrada en Sudamérica. No se identificaron
modelos entrenados para las especies chilenas amenazadas. La segunda es de
modalidad: la fusión de audio e imagen en fauna se ha explorado únicamente
con datos emparejados de especies norteamericanas, sin trabajos que la
apliquen a fuentes de ciencia ciudadana no emparejadas ni a avifaunas del
hemisferio sur. La tercera es de método: los SDM con aprendizaje profundo son
un área emergente que aún no se ha aplicado a la distribución de aves en
Chile, donde la práctica sigue apoyándose en MaxEnt. No existe, por tanto, un
sistema integrado que combine identificación multimodal y modelado de
distribución con aprendizaje profundo para la avifauna chilena. Esa es la
oportunidad que aborda este proyecto.
